% !PRINTED 334
\weberpage{334}
%<<CONT
set, where $f_{ik}$ stands for $f''(x_i,x_k)$. For the canonical form one obtains from this by a simple calculation
\weq{3}{\Delta(y) = -h\,m^2\,\Sigma y_1^3 + (h^3+2m^3)\,y_1y_2y_3.}
From this we derive the second covariant $C_1$ (\S.\,87), into which we likewise introduce a numerical factor, and accordingly
\weq{4}{J(x) = \frac{1}{36}\left|\begin{array}{cccc} f_{11}, & f_{12}, & f_{13}, & \Delta_1 \\ f_{21}, & f_{22}, & f_{23}, & \Delta_2 \\ f_{31}, & f_{32}, & f_{33}, & \Delta_3 \\ \Delta_1, & \Delta_2, & \Delta_3, & 0 \end{array}\right|}
set, where $\Delta_i$ stands for $\Delta'(x_i)$. This covariant too must be represented in the canonical form, which presents no difficulty if one forms the determinant according to the formula Bd. I, \S.\,23, (12). One finds thus:
\weq{5}{\begin{aligned}
J(y) &= -h^2(h^3+8m^3)^2\,\Sigma y_2^3y_3^3\\
&\quad + 3m^2(5h^6+26h^3m^3-4m^6)\,y_1^2y_2^2y_3^2\\
&\quad + hm(2h^6+5h^3m^3+20m^6)\,y_1y_2y_3\,\Sigma y_1^3\\
&\quad + 9h^2m^6(\Sigma y_1^3)^2.
\end{aligned}}
Accordingly the three symmetric functions $y_1y_2y_3$, $\Sigma y_1^3$, $\Sigma y_2^3y_3^3$ can be expressed through the three forms $\varphi(y)$, $\Delta(y)$, $J(y)$, and hence also through the corresponding forms in the original variables $x$, which we denote briefly by $f,\Delta,J$, according to the invariant relations
\weq{6}{\varphi(y) = f,\quad \Delta(y) = r^2\Delta,\quad J(y) = r^6 J}
One obtains first from (1) and (3):
\weq{7}{(h^3+8m^3)\,y_1y_2y_3 = m^2 f + r^2\Delta,}
\weq{8}{h(h^3+8m^3)\,\Sigma y_1^3 = (h^3+2m^3)\,f - 6\,m\,r^2\Delta,}
and then from (5)
\weq{9}{\begin{aligned}
h^2(h^3+8m^3)^2\,\Sigma y_2^3y_3^3 &= (2h^3+m^3)\,m^3 f^2 + m(2h^3-5m^3)\,r^2 f\Delta\\
&\quad + 3m^2 r^4 \Delta^2 - r^6 J.
\end{aligned}}
Very simple now is the calculation of the third covariant $C_2$ (\S.\,87) for the canonical form.

We introduce here too a simplifying numerical factor, and set
\weq{10}{K = K(x) = \frac{1}{9}\left|\begin{array}{ccc} f_1, & \Delta_1, & J_1 \\ f_2, & \Delta_2, & J_2 \\ f_3, & \Delta_3, & J_3 \end{array}\right|,}
%CONT>>
